\section{联络}
\label{sec:riem-connections}

度量在每个切空间内部比较向量，却不能比较两个不同点处的向量。单位圆上，$p=(1,0)$处的切向量$u=(0,1)$与$q=(0,1)$处的切向量$w=(-1,0)$都表示“逆时针前进一单位”，直觉上方向相同；可是作为$\mathbb R^2$中的向量，$u-w=(1,1)\neq0$，而且$u$根本不在$T_q\mathbb S^1$中。一般流形甚至没有环境空间可供相减。要谈“向量场沿某个方向变化了多少”，就必须规定怎样对向量场求导。

这项规定称为联络。它不是光滑结构自带的，同一个流形上可以有许多联络；有了黎曼度量之后，再加上两条自然的要求，就能从中唯一地选出一个。本节先给出联络的一般定义及其坐标计算，再确定后文一直使用的Levi-Civita联络。下文用$\mathfrak X(M)$记$M$上光滑向量场的全体，用$C^\infty(M)$记光滑实值函数的全体；向量场$X$作用于函数$f$得到方向导数$Xf=df(X)$，两个向量场的Lie括号$[X,Y]$满足$[X,Y]f=X(Yf)-Y(Xf)$，均见第~\ref{chap:manifold-basics}章。

\subsection{协变导数}
\label{sec:riem-covariant-derivative}

先看欧氏空间中我们想推广的东西。在$\mathbb R^n$中，向量场$Y$沿方向$X$的导数是对$Y$的每个直角坐标分量求方向导数：$D_XY=\sum_k(XY^k)\bm e_k$。它对$X$按函数线性（方向放大一倍，导数也放大一倍），对$Y$满足乘积法则：$D_X(fY)=(Xf)Y+fD_XY$。下面的定义只保留这些代数规则，而不再要求存在全局固定的基。

\begin{definition}[联络与协变导数]
\label{def:riem-connection}
光滑流形$M$上的\term{联络}（connection）是一个映射$\nabla:\mathfrak X(M)\times\mathfrak X(M)\to\mathfrak X(M)$，$(X,Y)\mapsto\nabla_XY$，对任意$X,X',Y,Y'\in\mathfrak X(M)$、$f,h\in C^\infty(M)$及$a,b\in\mathbb R$满足：
\begin{enumerate}
\item 对方向按函数线性：$\nabla_{fX+hX'}Y=f\nabla_XY+h\nabla_{X'}Y$；
\item 对被求导的场按实数线性：$\nabla_X(aY+bY')=a\nabla_XY+b\nabla_XY'$；
\item 乘积法则：$\nabla_X(fY)=(Xf)\,Y+f\,\nabla_XY$。
\end{enumerate}
$\nabla_XY$称为$Y$沿$X$的\term{协变导数}（covariant derivative）。
\end{definition}

第一条比第二条强：方向前面可以乘函数，被求导的场前面只能乘常数。原因在第三条：$Y$乘上函数$f$后，$f$本身的变化必须计入导数。第一条还带来一个重要推论：$(\nabla_XY)(p)$只依赖$X$在点$p$的值$X_p$，而不依赖$X$在$p$附近怎样取值。理由是在坐标中写$X=\sum_iX^i\partial_i$，由第一条，$\nabla_XY=\sum_iX^i\nabla_{\partial_i}Y$，在$p$点只出现数值$X^i(p)$。因此对单个切向量$v\in T_pM$也可以写$\nabla_vY$。对$Y$则没有这种性质，$\nabla_XY$要用到$Y$在$p$附近的变化。

\paragraph{局部系数}
在坐标图中，联络由$n^3$个函数完全确定。把坐标切基场代入，展开成同一组基的组合：
\begin{equation}
\nabla_{\partial_i}\partial_j=\sum_{k=1}^n\Gamma^k_{i,j}\,\partial_k .
\label{eq:riem-connection-coefficients}
\end{equation}
函数$\Gamma^k_{i,j}$称为联络在这张图中的系数；它们记录坐标基场沿坐标方向的变化。对一般的$X=\sum_iX^i\partial_i$与$Y=\sum_jY^j\partial_j$，由三条规则展开得
\begin{equation}
\nabla_XY=\sum_k\Bigl(\sum_iX^i\,\partial_iY^k+\sum_{i,j}\Gamma^k_{i,j}X^iY^j\Bigr)\partial_k .
\label{eq:riem-covariant-derivative-coordinates}
\end{equation}
第一部分是分量的普通导数，第二部分修正基场本身的变化。只有两部分合在一起才与坐标无关；单看分量导数，换一组坐标就会得到不同的结论。

\begin{example}[极坐标中的常向量场]
\label{ex:riem-polar-constant-field}
在欧氏平面上取通常的方向导数作联络，它在直角坐标中的系数全为零。改用极坐标$(r,\theta)$，由$\partial_r=(\cos\theta,\sin\theta)^{\top}$、$\partial_\theta=(-r\sin\theta,r\cos\theta)^{\top}$直接求导：
\begin{equation*}
\begin{gathered}
\nabla_{\partial_r}\partial_r=0,\qquad
\nabla_{\partial_r}\partial_\theta=\nabla_{\partial_\theta}\partial_r=\tfrac1r\partial_\theta,\\
\nabla_{\partial_\theta}\partial_\theta=(-r\cos\theta,-r\sin\theta)^{\top}=-r\,\partial_r .
\end{gathered}
\end{equation*}
所以非零系数只有$\Gamma^\theta_{r,\theta}=\Gamma^\theta_{\theta,r}=1/r$与$\Gamma^r_{\theta,\theta}=-r$。

考察常向量场$Y=\bm e_1$，即$Y=\cos\theta\,\partial_r-\frac{\sin\theta}r\partial_\theta$。它的分量随$\theta$变化，例如$\partial_\theta Y^r=-\sin\theta\neq0$，单看分量会误判$Y$在转动。用式\eqref{eq:riem-covariant-derivative-coordinates}计算$\nabla_{\partial_\theta}Y$：$r$分量为$-\sin\theta+\Gamma^r_{\theta,\theta}Y^\theta=-\sin\theta+(-r)(-\frac{\sin\theta}r)=0$，$\theta$分量为$-\frac{\cos\theta}r+\Gamma^\theta_{\theta,r}Y^r=-\frac{\cos\theta}r+\frac{\cos\theta}r=0$。同样可验证$\nabla_{\partial_r}Y=0$。修正项恰好抵消了基场的转动。
\end{example}

这个例子也说明，联络系数不是几何不变量：同一个联络在直角坐标中系数全为零，在极坐标中却不为零。因此不能凭系数是否为零判断几何性质，它们只是计算用的局部记录。

\paragraph{联络是额外的选择}
光滑结构并不决定联络。设$\nabla$是一个联络，$B(X,Y)$是任意一个对两个变量都按函数线性的向量值映射，例如在$\mathbb R^2$上取$B(X,Y)=X^1Y^1\bm e_2$，则$\nabla'_XY=\nabla_XY+B(X,Y)$仍满足定义中的三条规则：$B$对$Y$按函数线性，不会破坏乘积法则。在$\nabla'$下，常向量场$\bm e_1$沿$\bm e_1$方向的导数是$\bm e_2$，不再为零。反过来，任意两个联络之差都对两个变量按函数线性。哪一个联络合适，需要额外的要求来决定。

\paragraph{沿曲线的协变导数}
实际比较向量时，常常只有一条曲线和沿曲线给定的一族向量，例如飞机沿航线每时刻的速度。沿光滑曲线$\gamma:[a,b]\to M$的\emph{向量场}是指对每个$t$指定$V(t)\in T_{\gamma(t)}M$，且在坐标中$V(t)=\sum_jV^j(t)\partial_j$的分量光滑；它不必能延拓成$M$上的向量场，例如自交曲线在交点处可以有两个不同的值。对这样的$V$，联络给出沿曲线的协变导数
\begin{equation}
\frac{DV}{\mathrm dt}=\sum_k\Bigl(\dot V^k(t)+\sum_{i,j}\Gamma^k_{i,j}\bigl(x(t)\bigr)\,\dot x^i(t)\,V^j(t)\Bigr)\partial_k ,
\label{eq:riem-covariant-derivative-along-curve}
\end{equation}
其中$x(t)$是$\gamma$的坐标表示。它是式\eqref{eq:riem-covariant-derivative-coordinates}取$X=\dot\gamma$的结果：若$V$恰是某个向量场$Y$沿$\gamma$的取值，则$DV/\mathrm dt=\nabla_{\dot\gamma}Y$。可以验证这个表达式在换图时与坐标无关，并满足$\frac{D}{\mathrm dt}(fV)=\dot fV+f\frac{DV}{\mathrm dt}$。

\subsection{Levi-Civita联络}
\label{sec:riem-levi-civita}

在黎曼流形上，我们希望联络与度量配合。第一条要求是：对内积求导时，乘积法则照常成立，这样“两个向量都不变时它们的内积也不变”才有依据。第二条要求是：交换两个方向求导时，除了方向场本身不可交换造成的差（即Lie括号）以外，不再产生额外的差。欧氏空间的普通方向导数两条都满足。

\begin{definition}[度量相容与无挠]
\label{def:riem-compatible-torsion-free}
设$(M,g)$是黎曼流形，$\nabla$是$M$上的联络。若对任意$X,Y,Z\in\mathfrak X(M)$有
\begin{equation}
X\bigl(g(Y,Z)\bigr)=g(\nabla_XY,Z)+g(Y,\nabla_XZ),
\label{eq:riem-metric-compatible}
\end{equation}
称$\nabla$与$g$\term{度量相容}（metric compatibility）。若对任意$X,Y$有
\begin{equation}
\nabla_XY-\nabla_YX=[X,Y],
\label{eq:riem-torsion-free}
\end{equation}
称$\nabla$\term{无挠}（torsion-free）。同时满足两条的联络称为$g$的\term{Levi-Civita联络}（Levi-Civita connection）。
\end{definition}

两条要求在坐标中有直观的形式。坐标基场两两可交换，$[\partial_i,\partial_j]=0$，所以无挠等价于系数关于下标对称：$\Gamma^k_{i,j}=\Gamma^k_{j,i}$。把$\partial_i,\partial_j,\partial_\ell$代入度量相容条件，记$\Gamma_{i,j;\ell}=g(\nabla_{\partial_i}\partial_j,\partial_\ell)=\sum_k\Gamma^k_{i,j}g_{k,\ell}$，得到
\begin{equation}
\partial_ig_{j,\ell}=\Gamma_{i,j;\ell}+\Gamma_{i,\ell;j}.
\label{eq:riem-compatible-coordinates}
\end{equation}
也就是说，度量矩阵的变化完全由联络系数解释。注意无挠并不说$\nabla_X\nabla_YZ=\nabla_Y\nabla_XZ$；二阶导数的次序差是第~\ref{sec:riem-curvature}节的曲率。

\begin{theorem}[Levi-Civita联络的存在唯一性]
\label{thm:riem-levi-civita}
每个黎曼流形$(M,g)$上存在唯一的Levi-Civita联络。在任一坐标图中，它的系数为
\begin{equation}
\Gamma^k_{i,j}=\frac12\sum_{\ell=1}^ng^{k,\ell}\bigl(\partial_ig_{j,\ell}+\partial_jg_{i,\ell}-\partial_\ell g_{i,j}\bigr),
\label{eq:riem-christoffel}
\end{equation}
其中$[g^{k,\ell}]=\bm G^{-1}$。这些系数称为\term{Christoffel符号}（Christoffel symbols）。
\end{theorem}

\begin{proof}
先证唯一性。设$\nabla$是Levi-Civita联络。在式\eqref{eq:riem-compatible-coordinates}中轮换下标，写出$\partial_ig_{j,\ell}$、$\partial_jg_{i,\ell}$、$\partial_\ell g_{i,j}$三式，取前两式之和减去第三式：
\[
\partial_ig_{j,\ell}+\partial_jg_{i,\ell}-\partial_\ell g_{i,j}
=(\Gamma_{i,j;\ell}+\Gamma_{i,\ell;j})+(\Gamma_{j,i;\ell}+\Gamma_{j,\ell;i})-(\Gamma_{\ell,i;j}+\Gamma_{\ell,j;i}).
\]
无挠给出$\Gamma_{i,\ell;j}=\Gamma_{\ell,i;j}$、$\Gamma_{j,\ell;i}=\Gamma_{\ell,j;i}$与$\Gamma_{j,i;\ell}=\Gamma_{i,j;\ell}$，右端化为$2\Gamma_{i,j;\ell}$。再乘$\bm G^{-1}$解出$\Gamma^k_{i,j}$，即式\eqref{eq:riem-christoffel}。系数由度量唯一确定，而联络由系数唯一确定，所以在每张图上至多有一个Levi-Civita联络。

再证存在性。在一张坐标图中，用式\eqref{eq:riem-christoffel}定义系数，再用式\eqref{eq:riem-covariant-derivative-coordinates}定义$\nabla$，三条联络规则可直接验证。系数显然关于$i,j$对称，所以无挠。对度量相容，直接计算
\[
\Gamma_{i,j;\ell}+\Gamma_{i,\ell;j}=\tfrac12\bigl(\partial_ig_{j,\ell}+\partial_jg_{i,\ell}-\partial_\ell g_{i,j}\bigr)+\tfrac12\bigl(\partial_ig_{\ell,j}+\partial_\ell g_{i,j}-\partial_jg_{i,\ell}\bigr)=\partial_ig_{j,\ell},
\]
即式\eqref{eq:riem-compatible-coordinates}对坐标基场成立；一般向量场的情形由两边对$Y,Z$的乘积法则展开后归结为这一情形。最后，两张图的重叠处，两种定义都是同一个度量的Levi-Civita联络，由唯一性它们一致，于是各图的定义拼成整体的联络。
\end{proof}

证明中的轮换组合也可以不借坐标写出，称为Koszul公式；它表明$g(\nabla_XY,Z)$可以只用度量、方向导数和Lie括号表示。定理的要点是：一旦给定度量，就不再有选择联络的自由，后文所说的联络、平行移动、测地线和曲率，若不另作说明，都指Levi-Civita联络给出的对象。

\begin{example}[由度量求Christoffel符号]
\label{ex:riem-christoffel-examples}
欧氏平面的极坐标中$\bm G=\operatorname{diag}(1,r^2)$，只有$g_{\theta,\theta}=r^2$依赖坐标。代入式\eqref{eq:riem-christoffel}：
\begin{equation*}
\begin{aligned}
\Gamma^r_{\theta,\theta}&=\frac12g^{r,r}(-\partial_rg_{\theta,\theta})=-r,\\
\Gamma^\theta_{r,\theta}&=\frac12g^{\theta,\theta}\partial_rg_{\theta,\theta}=\frac1{2r^2}\cdot2r=\frac1r,
\end{aligned}
\end{equation*}
其余为零，与例~\ref{ex:riem-polar-constant-field}直接求导的结果一致。

单位球面在例~\ref{ex:riem-sphere-metric}的坐标中$\bm G=\operatorname{diag}(1,\sin^2\theta)$，同样计算得到非零系数
\[
\Gamma^\theta_{\phi,\phi}=-\sin\theta\cos\theta,\qquad\Gamma^\phi_{\theta,\phi}=\Gamma^\phi_{\phi,\theta}=\cot\theta .
\]
这两组系数在后面的测地线和曲率计算中还会用到。
\end{example}

\paragraph{嵌入子流形上的投影公式}
对取诱导度量的嵌入子流形$M\subseteq\mathbb R^D$，Levi-Civita联络有一个不用坐标的描述。把切向量场$Y$看成取值于$\mathbb R^D$的函数，先求它沿$X$的普通方向导数$D_XY$，再用正交投影$\bm P_p:\mathbb R^D\to T_pM$去掉法向部分：
\begin{equation}
(\nabla_XY)(p)=\bm P_p\bigl((D_XY)(p)\bigr).
\label{eq:riem-projection-connection}
\end{equation}
这确实是Levi-Civita联络：三条联络规则由普通导数继承；对切向量$Y,Z$，$X\langle Y,Z\rangle=\langle D_XY,Z\rangle+\langle Y,D_XZ\rangle$，而$Z$切向，内积只读取$D_XY$的切向部分，所以度量相容；$D_XY-D_YX=[X,Y]$本身已是切向量，投影后不变，所以无挠。由唯一性，它就是诱导度量的Levi-Civita联络。在单位圆上，匀速逆时针的单位切向量场环境导数指向圆心，投影后为零：按联络判断，这个场沿圆周没有变化，正符合开头的直觉。
